% --- Άσκηση 14811 (14811.tex) ---
\Askhsh{14811}{1}
\begin{ylh}{2}
    \item 2.3. Απόλυτη Τιμή Πραγματικού Αριθμού
\end{ylh}
Να χαρακτηρίσετε καθεμιά από τις προτάσεις που ακολουθούν ως \textbf{Σωστή (Σ)} ή \textbf{Λανθασμένη (Λ)}, γράφοντας στην κόλλα σας, δίπλα στο αριθμό που αντιστοιχεί σε καθεμιά από αυτές το γράμμα Σ αν η πρόταση είναι Σωστή, ή το γράμμα Λ αν αυτή είναι Λάθος.
\begin{alist}
\item
  Το σημείο $M(x,y)$ με $x > 0$ και $y < 0$ βρίσκεται στο δεύτερο τεταρτημόριο του καρτεσιανού συστήματος συντεταγμένων.
\item
  Αν τρεις μη μηδενικοί αριθμοί $a,\beta,\gamma$ είναι διαδοχικοί όροι γεωμετρικής προόδου, τότε ισχύει: $\beta^{2} = a \cdot \gamma$.
\item
  Ισχύει $|a| \geq a$, για κάθε $a \in \mathbb{R}$.
\item
  Αν $a > \beta$ και $\gamma > \delta$, τότε: $a - \gamma > \beta - \delta$ για οποιουσδήποτε πραγματικούς αριθμούς α, β, γ, δ.
\item
  Η εξίσωση $a x = a$ έχει μοναδική λύση $x = 1$ για κάθε $a\mathbb{\in R}$ .
\monades{10}
\item Για τους πραγματικούς αριθμούς $a,\beta$ να αποδείξετε ότι: $|a \cdot \beta|\  = |a| \cdot |\beta|$.
\monades{15}
\end{alist}
